% =========================================
% PHẦN 1: Đặt vấn đề và Mục tiêu nghiên cứu
% =========================================

\begin{frame}
  \centering
  \Huge \textbf{PHẦN 1} \\[0.5cm]
  \LARGE Đặt vấn đề và Mục tiêu nghiên cứu
\end{frame}

%---------------------------------
\begin{frame}{Bài toán điều khiển đèn giao thông}
\begin{itemize}
    \item \textbf{Bối cảnh:} Ùn tắc giao thông đô thị ngày càng nghiêm trọng tại Hà Nội và TP.HCM, gây thiệt hại kinh tế hàng tỷ đồng mỗi ngày.
    \item \textbf{Đặc thù giao thông Việt Nam:}
    \begin{itemize}
        \item Xe máy chiếm $\approx$60--70\% tổng lưu lượng — cao hơn nhiều so với các nước phát triển
        \item Hệ số PCU: 1 ô tô $\approx$ 3 xe máy về chiếm dụng diện tích
        \item Lưu lượng bất đối xứng giữa các trục đường (ví dụ: Đông–Tây gấp đôi Bắc–Nam)
    \end{itemize}
    \item \textbf{Hạn chế phương pháp truyền thống:}
    \begin{itemize}
        \item Điều khiển thời gian cố định (Fixed-Time): chu kỳ cứng nhắc, không thích ứng
        \item Không phản ứng với biến động lưu lượng thực tế
        \item Hầu hết nghiên cứu hiện có xây dựng cho giao thông Châu Âu/Bắc Mỹ
    \end{itemize}
    \item \textbf{Mục tiêu điều khiển:} Phân bổ thời gian xanh và chuyển pha tối ưu nhằm giảm ùn tắc, thích ứng theo lưu lượng thực tế.
\end{itemize}
\end{frame}

%---------------------------------
\begin{frame}{Giao thông như một bài toán ra quyết định tuần tự}
\begin{itemize}
    \item \textbf{Định nghĩa:} Điều khiển đèn giao thông diễn ra theo chuỗi thời gian, quyết định hiện tại ảnh hưởng đến trạng thái tương lai.
    \item \textbf{Đặc điểm:}
    \begin{itemize}
        \item \textbf{Sự đánh đổi:} Ưu tiên một hướng có thể gây ùn tắc ở hướng khác.
        \item \textbf{Tính động:} Hệ thống tiến triển liên tục; trạng thái kế tiếp phụ thuộc trạng thái hiện tại và hành động đã thực hiện.
        \item \textbf{Tính ngẫu nhiên:} Phương tiện đến ngẫu nhiên, không thể dự đoán chính xác.
    \end{itemize}
    \item \textbf{Yêu cầu:} Cần một khung mô hình hóa được sự tiến triển trạng thái và tác động dài hạn của quyết định.
    \item \textbf{Hướng tiếp cận:} Mô hình hóa bài toán dưới dạng \textbf{MDP} và áp dụng \textbf{Deep Reinforcement Learning (Double DQN)} để học chính sách điều khiển tối ưu.
\end{itemize}
\end{frame}

%---------------------------------
\begin{frame}{Mục tiêu dự án}
\begin{enumerate}
    \item \textbf{Xây dựng môi trường mô phỏng} ngã tư Hà Nội mẫu bằng \textbf{SUMO} — tích hợp phân bố xe đặc trưng VN (60\% xe máy) và hệ số PCU chuẩn hóa.

    \vspace{4pt}
    \item \textbf{Mô hình hóa bài toán} dưới dạng Mô hình Quyết định Markov (MDP): xác định không gian trạng thái (32 chiều), không gian hành động (4 pha đèn) và hàm phần thưởng phản ánh mức độ ùn tắc.

    \vspace{4pt}
    \item \textbf{Triển khai thuật toán Double DQN kết hợp Dueling Network} để huấn luyện agent tự thích ứng với biến động giao thông theo thời gian thực.

    \vspace{4pt}
    \item \textbf{Đánh giá định lượng} DQN agent so với baseline Fixed-Time 160s qua các chỉ số: hàng đợi, thời gian chờ, tốc độ lưu thông và lưu lượng xe qua nút.
\end{enumerate}
\end{frame}
